\chapter{Conclusion}
\label{chap:conclusion}


% =================================================================
% 6.1 Conclusions
% =================================================================

\section{Conclusions}
\label{sec:conclusion-summary}

This thesis investigated whether self-supervised vision foundation
models can be adapted for precise fetal biometry landmark
localisation. To bridge the gap between general-purpose visual
representations and the spatially precise, measurement-specific
predictions required for fetal biometry, it introduced QCLand, a
query-conditioned landmark-localisation framework that combines
learned landmark queries with pretrained DINOv2 and DINOv3 backbones
and a new spatial decoding head, DeconvHeadV2.

The final comparison provides evidence that QCLand can be adapted
effectively to this task. A QCLand variant records the lowest mean
PI-NME in eight of the ten dataset--measurement settings, with
QCLand-DINOv3 achieving the lowest mean on all five UCL
measurements. These results support query-conditioned
foundation-model adaptation as a viable approach to precise fetal
landmark localisation within the evaluated benchmark settings.

The BPD ablation study further clarifies where the observed
performance gains arise. DeconvHeadV2 produces the largest reduction
in localisation error among the evaluated architectural
modifications. Rotation-and-scale augmentation is associated with a
further substantial reduction in mean error and lower run-to-run
variability. By contrast, the independent contributions of
multi-depth feature fusion and pixel-centre alignment remain
unresolved by the sequential experimental design.

Performance nevertheless varies across anatomy, backbone, and dataset
setting. DINOv3 is descriptively stronger than DINOv2 in most
comparisons, but paired statistical support is confined to two
Multicentre measurements. The QCLand--HRNet-W18 ranking is similarly
anatomy-dependent: QCLand shows its clearest confirmed advantage on
Multicentre APAD, whereas HRNet-W18 retains a clear advantage on the
Multicentre head measurements. RTMPose-s is consistently weaker on
UCL but remains competitive in several Multicentre settings.

QCLand is therefore competitive across the evaluated fetal
landmark-localisation tasks, but its advantages are conditional rather
than universal. The results support the effectiveness of
landmark-query conditioning and spatial decoding while also showing
that architecture, backbone, anatomy, and dataset characteristics
remain important determinants of performance.


% =================================================================
% 6.2 Contributions
% =================================================================

\section{Contributions}
\label{sec:conclusion-contributions}

The principal methodological contribution of this thesis is QCLand,
a framework for adapting pretrained vision foundation models to
precise fetal landmark localisation. QCLand introduces learned
landmark queries into the pretrained transformer, allowing
landmark-specific representations to emerge through direct
query--image interaction, and couples these representations with
DeconvHeadV2 for explicit spatial decoding. This provides a route
from general-purpose transformer representations to precise
landmark-level predictions.

A second contribution is the systematic development of this
adaptation strategy. The BPD ablation study traces the transition
from the inherited dot-product readout to the adopted QCLand
configuration and identifies DeconvHeadV2 as the architectural change
associated with the largest observed improvement. The experiments
also characterise the effects and remaining uncertainty associated
with the training objective, multi-depth feature fusion,
pixel-centre alignment, geometric augmentation, and input resolution.

The empirical contribution is a controlled comparison across five
fetal measurements, two dataset settings, two foundation-model
backbones, and two established landmark-localisation baselines.
QCLand-DINOv2 and QCLand-DINOv3 are evaluated against locally
reproduced HRNet-W18 and RTMPose-s models under harmonised external
conditions using five training seeds, common evaluation subsets,
original-image-space PI-NME, and paired statistical analysis. This
evaluation shows both where QCLand is competitive and where
purpose-built alternatives retain an advantage.


% =================================================================
% 6.3 Limitations and Future Work
% =================================================================

\section{Limitations and Future Work}
\label{sec:conclusion-future}

Several limitations constrain the conclusions that can be drawn from
the present experiments. The Multicentre evaluation uses the released
pooled benchmark rather than a centre-held-out design, so it does not
directly test transfer to an unseen clinical site. UCL is also one of
the constituent sources of the Multicentre release, meaning that the
two dataset settings should not be interpreted as fully independent
replications.

The compared methods retain different training procedures despite the
harmonised external evaluation protocol. QCLand uses
validation-based termination, HRNet-W18 follows its upstream training
procedure, and RTMPose-s trains for a fixed number of epochs.
Consequently, the comparison evaluates complete methods under matched
external conditions rather than isolating architecture alone.

The BPD ablation study is staged rather than factorial. The observed
effect of each modification therefore depends on the configuration
that precedes it, and multi-depth fusion and pixel-centre alignment
are varied together in the augmented follow-up. Their independent
contributions cannot be established from the present evidence.
Furthermore, the full ablation sequence was conducted only on BPD,
while the remaining measurements evaluate the adopted configuration
as a complete pipeline.

Training-time landmark assignment provides another direction for
further investigation. QCLand assigns the two landmarks from left to
right, and this convention can become unstable when their horizontal
separation is small. PI-NME removes the corresponding ordering
ambiguity during evaluation, but it does not alter the query-specific
supervision used during training. A per-image analysis relating
landmark orientation to localisation error would test whether this
mechanism contributes to the observed anatomy-dependent performance.
If supported, orientation-aware canonicalisation or
permutation-aware training objectives would provide natural
extensions.

Further experiments could also clarify the effects of dataset scale
and anatomical variation. A controlled subsampling analysis on
Multicentre would help distinguish the influence of training-set size
from source and device diversity, while extending the ablation study
to a non-head measurement would test whether the strong
DeconvHeadV2 result and the less conclusive fusion/alignment effects
generalise beyond BPD.

The implementation and reproduction scripts will be made publicly
available at
\url{https://github.com/Flora020703/qcland-landmark-detection}
upon submission.


% =================================================================
% 6.4 Final Perspective
% =================================================================

\section{Final Perspective}
\label{sec:conclusion-closing}

The results of this thesis suggest that strong pretrained visual
representations alone are not sufficient for precise structured
localisation. Effective adaptation also requires mechanisms that
preserve spatial precision while allowing the representation to
specialise to the structure of the prediction task. QCLand addresses
this requirement by combining landmark-query interaction with
query-conditioned spatial decoding.

Across the evaluated fetal biometry tasks, this design is competitive
with established convolutional and coordinate-classification
approaches while exhibiting clear anatomy- and dataset-dependent
behaviour. The broader implication is therefore not that foundation
models universally replace purpose-built localisation architectures,
but that their representations can become effective for precise
anatomical landmark localisation when paired with an appropriate
task-specific adaptation mechanism.
